\chapter*{Abstract}
\addcontentsline{toc}{chapter}{Abstract}

This thesis investigates how reliability, system stability, and execution performance change
when a workflow orchestration system is migrated from a single virtual machine (\vm{}) using
a Docker executor (where each pipeline job runs as an isolated Docker container on the VM
host) to a local Kubernetes cluster using the Kubernetes Run Launcher with isolated pod
execution. \dagster{} is used as the representative workflow orchestration framework. The
study is conducted entirely on local infrastructure: a UTM Ubuntu 22.04 \textsc{vm}
provisioned via Ansible for the baseline condition and a single-node Kind cluster on the
same host for the Kubernetes condition, ensuring a fair resource comparison and eliminating
cloud billing constraints.

The central research problem is that engineering teams routinely make \vm{}-to-Kubernetes
migration decisions based on industry assumptions rather than measured behavioural evidence.
The primary contribution of this thesis is the identification and empirical measurement of a
\textit{crossover point} --- formally defined as the concurrent workload level at which
the \vm{} deployment's job success rate drops below 95\% and the Kubernetes deployment's
total execution time including scheduling overhead becomes lower than the \vm{}'s execution
time under active resource contention.

Four research questions are investigated through three controlled experiments measuring
job success rate, execution time variance, failure blast radius, Kubernetes scheduling
overhead, and resource utilisation across six defined workload levels (1, 2, 4, 8, 12, and
16 concurrent jobs). Each experiment is repeated three times per level and results are
reported with standard deviation to ensure statistical credibility.

The study is grounded in four literature pillars: resource contention and performance
collapse theory, Kubernetes overhead research, autoscaling response behaviour, and
distributed systems fault isolation.

\textbf{Keywords:} Kubernetes, Kind, Dagster, workflow orchestration, resource contention,
pod isolation, crossover point, reliability measurement, DevOps

\newpage
